\section{The motivic dictionary used in the proof}
\label{cp:sec-motivic-dictionary}

This section is a dictionary, not a second proof.  It explains the objects
which occur in the motivic step and records exactly which features of the
general theory are used.  A reader may take the properties listed below as
axioms for the argument.  In particular, we do not require a conjectural
abelian category of all mixed motives.

\subsection{From a spherical simplex to an oriented quadric simplex}
\label{cp:subsec-quadric-simplex}

Let \(U\) be a smooth connected variety over \(\Q\), and let \(E\) be a
rank-four vector bundle on \(U\).  A nowhere-degenerate quadratic form on
\(E\) cuts out a smooth relative quadric
\[
 Q\subset \mathbf P(E).
\]
Let \(M=(M_0,M_1,M_2,M_3)\) be four ordered relative hyperplanes in general
position.  We require the intersections of \(Q\) with all coordinate faces
to have the expected dimension and to be regular.  The pair \((Q,M)\) is the
algebro-geometric version of a tetrahedron: the quadric is the complexified
null cone which encodes the spherical metric (classically called the
\emph{absolute}), and the \(M_i\) are its four faces.

In the family constructed in Section~\ref{cp:sec-phase-charts} this data is
concrete:
\[
 E=\mathcal O_U^4,\qquad M_i=\{z_i=0\}.
\]
The spherical simplex gives the quadratic form with matrix \(G^{-1}\), hence
\(Q_{G^{-1}}\); projective duality replaces it by \(Q_G\), on which the
coordinate-face calculation is easier.  The function \(w\), satisfying
\(w^2=\det G\), chooses one of the two rulings and therefore orients both
descriptions.

For a quadric surface, the maximal isotropic subspaces are lines.  Over an
algebraic closure they occur in two rulings.  An \emph{orientation} of \(Q\)
is a choice of one of these two components, made compatibly in the family.
For a matrix presentation
\[
 Q_G:\quad z^tGz=0,
\]
the two choices are distinguished by a square root \(w\) of \(\det G\).
Changing \(w\) interchanges the rulings and reverses the oriented
quadric-simplex class.  Thus the natural parameter space for oriented
quadrics is the determinant double cover \(w^2=\det G\), rather than the
space of matrices \(G\) alone.

The notation
\[
 [(Q,\alpha);M_0,M_1,M_2,M_3]
\]
means the resulting oriented quadric simplex, where \(\alpha\) is the chosen
ruling.  Goncharov's group \(S_4(F)\), for a field \(F\), is generated by
such objects over \(F\), modulo projective invariance, alternating face
relations, and the standard additivity relations.  We will not need a
presentation of all these relations; we use only its face coproduct and the
canonical maps listed in Section~\ref{cp:subsec-standard-inputs}.

There are two equivalent matrix descriptions which occur in the proof.
Start with
\[
 Q_{G^{-1}}:\quad x^tG^{-1}x=0\subset\mathbf P(E)
\]
and the coordinate simplex.  Goncharov's projective-duality operation
replaces the face opposite the vertex \(e_j\) by its polar hyperplane
\[
 e_j^tG^{-1}x=0.
\]
Now make the projective change of coordinates \(x=Gy\).  The quadric equation
becomes \(y^tGy=0\), and the four polar hyperplanes become \(y_j=0\).
Thus the dual class is represented by \(Q_G\) with the coordinate faces.
This operation is an involutive linear automorphism of quadric-scissors
classes, so a class is zero if and only if its dual is zero.  This is why the
face calculation may be performed on \(Q_G\), even though the quadric first
attached to the Gram matrix is \(Q_{G^{-1}}\).  The determinant root
transports the orientation through the operation.

\subsection{What a relative motive means here}
\label{cp:subsec-what-is-a-motive}

A cohomology theory replaces a variety by groups.  A motive retains the
common algebraic structure which maps to all the usual realizations---for
example Betti, de Rham, or Hodge cohomology.  In this proof it is important
that a motive is an \emph{object}, not merely a numerical invariant: it can
be put into exact triangles, dualized, filtered, and pulled back to a point.

We work in the rational constructible motivic category
\[
 \operatorname{DM}(U)_\Q.
\]
For the present argument this may be understood as a symmetric monoidal,
idempotent-complete triangulated category with the following operations.
Our cohomological shift convention is
\(H^i(A[m])=H^{i+m}(A)\) in every realization.  A distinguished triangle
\(A\to B\to C\to A[1]\) is the derived replacement for a short exact
sequence: every cohomology realization turns it into a long exact
sequence.  Thus the triangles below are not extra geometric objects; they
are the bookkeeping device which remembers kernels, quotients, and
connecting maps simultaneously.

\begin{enumerate}[label=(\roman*)]
\item It has a unit object \(\Q_U\), shifts \(A[m]\), tensor products, and
      internal Homs \(\underline{\operatorname{Hom}}(A,B)\).

\item A smooth \(U\)-scheme \(f:X\to U\) has a geometric motive \(M_U(X)\).
      Disjoint unions become direct sums and \(M_U(U)=\Q_U\).

\item If \(Z\hookrightarrow X\) is a regular closed immersion with open
      complement \(V\), localization gives a distinguished triangle.  With
      purity, for codimension \(c\) it can be written in Gysin form
      \[
       M_U(V)\longrightarrow M_U(X)\longrightarrow
       M_U(Z)(c)[2c]\longrightarrow M_U(V)[1].
      \]
      Repeated localization is the motivic version of cutting a space into
      strata.

\item Constructible objects in the calculation are dualizable.  We write
      \[
       D(A)=\underline{\operatorname{Hom}}(A,\Q_U).
      \]
      Poincar\'e duality relates this internal dual to compactly supported
      cohomology in every realization.

\item For a geometric point \(s:\operatorname{Spec}k\to U\), exact symmetric
      monoidal pullback
      \[
       s^*: \operatorname{DM}(U)_\Q\longrightarrow
             \operatorname{DM}(k)_\Q
      \]
      commutes with every construction used below: geometric motives,
      localization maps, tensor products, twists, idempotent summands, and
      internal duals.

\item A finite \'etale map has pullback, pushforward, a unit, and a trace.
      For a degree-\(d\) cover, trace after the unit is multiplication by
      \(d\).  This is the only finite-\'etale descent fact used below.
\end{enumerate}

These are the only formal properties of the ambient motivic category that
enter the proof.

For later use, \emph{motivic cohomology} is simply the following Hom group in
this category:
\begin{equation}
 H^m(U,\Q(n))=
 \operatorname{Hom}_{\operatorname{DM}(U)_\Q}
  \bigl(\Q_U(0),\Q_U(n)[m]\bigr).
 \label{cp:eq-motivic-cohomology-definition}
\end{equation}
Thus every motivic-cohomology class used in the proof is an actual morphism
between two explicitly named objects.

\subsection{Tate objects and Artin objects}
\label{cp:subsec-tate-artin}

The \emph{Tate object} \(\Q_U(1)\) is the basic one-dimensional object
detected by the multiplicative group.  Our convention is
\[
 \Q(1)=\mathbf G_m[-1]
\]
where on the right \(\mathbf G_m\) denotes the motivic sheaf represented by
the multiplicative group; this is an equality of motivic objects, not of
schemes.
Concretely, on every smooth connected scheme \(W\) used here,
\begin{equation}
 H^1(W,\Q(1))\cong\mathcal O(W)^*\otimes_{\mathbf Z}\Q.
 \label{cp:eq-weight-one-units}
\end{equation}
This is the bridge between the angle and length cross-ratios, which are
invertible functions, and the middle Tate degree of the motive.
For any object \(A\), its \(n\)-fold Tate twist is
\[
 A(n)=A\otimes\Q_U(n),
 \qquad
 \Q_U(n)^\vee=\Q_U(-n).
\]
The integers \(0,1,2\) used below are \emph{Tate degrees}.  They are the
three levels relevant to a three-dimensional quadric simplex; they should
not be confused with the signs or factor-of-two conventions used for
Deligne's numerical weights.

An \emph{Artin motive} is the motivic avatar of a finite \'etale cover.  The
only nonsplit one needed here comes from a quadratic cover
\[
 \pi:U'\longrightarrow U
\]
with deck involution \(\sigma\).  Rationally,
\begin{equation}
 M_U(U')=\pi_*\Q_{U'}\cong \Q_U\oplus\chi,
 \qquad
 \chi=\operatorname{im}\frac{1-\sigma}{2}.
 \label{cp:eq-sign-motive}
\end{equation}
The first summand is invariant and the second is the sign motive.  More
intrinsically, if \(D\to U\) is any degree-two finite \'etale scheme, its
reduced motive is the idempotent summand
\begin{equation}
 \widetilde M(D)=
 \operatorname{im}\left(
  1-\frac12\eta\operatorname{Tr}:M_U(D)\longrightarrow M_U(D)
 \right),
 \label{cp:eq-reduced-artin-motive}
\end{equation}
where \(\eta\) is the unit and \(\operatorname{Tr}\) is the trace.  For a
connected quadratic cover this is \(\chi\).  For a split pair of sections it
is a copy of \(\Q_U\).  The rational Artin category is semisimple: for a
Galois cover this is simply Maschke's theorem for rational representations
of its finite deck group.

\subsection{The alternating face complex}
\label{cp:subsec-face-complex}

\begin{definition}[Normalized relative quadric motive]
\label{cp:def-face-motive}
For \(I\subset\{0,1,2,3\}\), put
\[
 M_I=\bigcap_{i\in I}M_i,
 \qquad Q_I=Q\cap M_I.
\]
The order of the hyperplanes supplies the usual alternating incidence
signs.  Place the \(|I|=3\) term in cohomological degree \(0\), and the
following three terms in degrees \(1,2,3\), respectively.  The geometric
face complex is the total object
\begin{equation}
\begin{split}
 \bigoplus_{|I|=3}M_U(M_I\setminus Q_I)
 &\longrightarrow
 \bigoplus_{|I|=2}M_U(M_I\setminus Q_I)\\
 &\longrightarrow
 \bigoplus_{|I|=1}M_U(M_I\setminus Q_I)
 \longrightarrow M_U(\mathbf P(E)\setminus Q).
\end{split}
 \label{cp:eq-geometric-face-complex}
\end{equation}
The phrase \emph{total object} means the finite totalization of this
four-term complex.  We use the standard stable (equivalently, dg)
enhancement of \(\operatorname{DM}(U)_\Q\), in which a bounded diagram has a
functorial totalization.  Its image in the triangulated homotopy category may
be computed by iterated mapping cones.
We denote it by \(\widetilde m_U(Q,M)\).  Its raw cohomological dual and its
normalized version are
\begin{equation}
\begin{aligned}
 h_U(Q,M)&=D\bigl(\widetilde m_U(Q,M)\bigr),\\
 \mathfrak m_U(Q,M)&=h_U(Q,M)(2)\\
 &=\underline{\operatorname{Hom}}
       \bigl(\widetilde m_U(Q,M),\Q_U\bigr)(2).
\end{aligned}
 \label{cp:eq-normalized-relative-motive}
\end{equation}
\end{definition}

The face complex is the motivic analogue of the relative cellular complex
of the complement of \(Q\), relative to its four ordered faces.  Its
alternating signs record how a face occurs in the boundary of the
next-dimensional face.  This description is useful because localization
can be applied to each displayed stratum separately.

The twist \((2)\) in Definition~\ref{cp:def-face-motive} is essential and
there is no additional cohomological shift.  Before twisting, the two
endpoints have Tate degrees \(0\) and \(-2\); after twisting they become
\(2\) and \(0\).  This is precisely the normalization in which a primitive
endpoint is an extension of \(\Q_U(0)\) by \(\Q_U(2)\).

\subsection{Why only three weights occur}
\label{cp:subsec-three-weights}

Assume \(Q\subset\mathbf P^3_U\) is oriented and split, every plane section
\(Q\cap M_i\) is a split regular conic, and every line section
\[
 D_I=Q\cap M_I\subset M_I\cong\mathbf P^1_U,
 \qquad |I|=2,
\]
is finite \'etale of degree two.  Localization, purity, and the projective
bundle formula give
\begin{align}
 M_U(\mathbf P^3_U\setminus Q)
   &\cong \Q_U\oplus\Q_U(2)[3],
 \label{cp:eq-open-quadric-motive}\\
 M_U(\mathbf P^2_U\setminus C)
   &\cong \Q_U
   &&(C\text{ a split regular conic}),
 \label{cp:eq-open-conic-motive}\\
 M_U(\mathbf P^1_U\setminus D_I)
   &\cong \Q_U\oplus\widetilde M(D_I)(1)[1],
 \label{cp:eq-punctured-line-motive}\\
 M_U(\mathbf P^0_U)&\cong\Q_U.
 \label{cp:eq-point-motive}
\end{align}
The shifts \([3]\) and \([1]\) cancel the positions of the corresponding
terms in the face complex.  After internal duality and the twist \((2)\), the
normalized relative motive therefore has a three-step filtration with
graded pieces
\begin{equation}
 \operatorname{gr}^0\mathfrak m_U\cong\Q_U(0),
 \qquad
 \operatorname{gr}^1\mathfrak m_U\cong
   \bigoplus_{|I|=2}\widetilde M(D_I)(1),
 \qquad
 \operatorname{gr}^2\mathfrak m_U\cong\Q_U(2).
 \label{cp:eq-dictionary-three-graded-pieces}
\end{equation}
We use the decreasing convention
\(F^0\supset F^1\supset F^2\supset F^3=0\), with
\(\operatorname{gr}_F^i=F^i/F^{i+1}\).  Thus \(\Q_U(2)\) is the bottom
subobject \(F^2\), while \(\Q_U(0)\) is the top quotient
\(F^0/F^1\).  This convention fixes the direction of the endpoint
extension used below.
This formula is the reason the argument is called three-weight.  All other
summands in the dual cellular row cancel, and the three displayed kinds of
objects are the complete list of its graded pieces.  If five binary sections
split and one is the quadratic cover \(U'\to U\), then the middle term is a
direct sum of copies of \(\Q_U(1)\) and one possible copy of \(\chi(1)\).

\subsection{The restricted heart and extensions}
\label{cp:subsec-restricted-heart}

One often packages mixed motives into an abelian category, the \emph{heart}
of a motivic \(t\)-structure.  A \(t\)-structure is a rule separating a
triangulated category into nonpositive and nonnegative parts; their
intersection is an abelian category, and triangles among its objects are
short exact sequences.  The existence of such a heart for every
motive over every base is not assumed here.  We use only the finite window
generated by
\begin{equation}
 \mathcal P_0=\langle\Q_U(0)\rangle,\qquad
 \mathcal P_1=\langle\Q_U(1),\chi(1)\rangle,\qquad
 \mathcal P_2=\langle\Q_U(2)\rangle.
 \label{cp:eq-pure-three-levels}
\end{equation}
Each \(\mathcal P_i\) is semisimple.  If the negative Hom groups between
these three pure levels vanish, the finite semisimple gluing theorem equips
their thick triangulated closure with a bounded \(t\)-structure.  Its heart,
denoted
\[
 \mathcal M_U^{[0,2]},
\]
is the extension closure of the three \(\mathcal P_i\).  Here ``thick'' means
closed under shifts, cones, finite sums, and direct summands; ``extension
closure'' means objects built from the displayed pure objects by finitely
many short exact sequences; and a negative Hom means
\(\operatorname{Hom}(P_i,P_j[m])\) with \(m<0\).  Concretely, its
objects are assembled from the three pure levels by finitely many short exact
sequences.  A distinguished triangle whose three terms lie in the heart is
the corresponding short exact sequence.  To obtain the functorial
three-step filtration one also uses the vanishing of wrong-way
degree-zero maps and degree-one extensions.  Any extension written in the
wrong weight order then splits and may be reordered; induction on extension
length gives a unique maximal subobject of weights at least \(i\).  These
subobjects are \(F^i\), and their uniqueness makes every morphism strict,
meaning \(f(F^iM)=\operatorname{im}(f)\cap F^iN\) for every morphism
\(f:M\to N\).
Section~\ref{cp:sec-relative-rigidity} verifies these additional vanishings
in the present motivic window.
Inside the restricted heart, the heart-to-derived comparison gives
\begin{equation}
 \operatorname{Ext}^1_{\mathcal M_U^{[0,2]}}
   \bigl(\Q_U(0),\Q_U(2)\bigr)
 \cong H^1(U,\Q(2)).
 \label{cp:eq-ext-is-motivic-cohomology}
\end{equation}
An element of this group is represented by an exact sequence
\[
 0\longrightarrow\Q_U(2)\longrightarrow E\longrightarrow
 \Q_U(0)\longrightarrow0,
\]
or equivalently by the endpoint morphism
\(\Q_U(0)\to\Q_U(2)[1]\) in its distinguished triangle.

\subsection{Frames, coproducts, and primitive classes}
\label{cp:subsec-frames-coproduct}

\begin{definition}[Framed class]
\label{cp:def-framed}
A \((0,2)\)-framing on an object of \(\mathcal M_U^{[0,2]}\) specifies its
two endpoints: a chosen identification of its degree-zero graded piece with
\(\Q_U(0)\), and a chosen identification of its degree-two graded piece with
\(\Q_U(2)\).  Equivalently, after applying the graded fibre functor defined
below, one gives the corresponding nonzero top vector and bottom covector.
Framed equivalence remembers these maps while
allowing changes to the unframed middle of the object.
\end{definition}

A mixed object alone forgets which geometric fundamental classes are being
compared.  For \(\mathfrak m_U(Q,M)\), the chosen ruling of the oriented
quadric supplies the top \(\Q_U(0)\)-vector, and the ordering and
orientation of the four coordinate faces supplies the bottom
\(\Q_U(2)\)-covector.  Thus an oriented quadric simplex
determines a framed class, not just an unmarked motive.

\begin{definition}[Coefficient spaces]
\label{cp:def-coefficient-space}
Let \(X\) be simple in \(\mathcal P_i\), let \(Y\) be simple in
\(\mathcal P_j\), and assume \(i<j\).  An \((X,Y)\)-framed coefficient is
a triple
\[
 (M,v,f),\qquad
 v:X\longrightarrow\operatorname{gr}_F^iM,\qquad
 f:\operatorname{gr}_F^jM\longrightarrow Y.
\]
Thus \(v\) is the top vector and \(f\) the bottom covector.  Two triples are
identified when a morphism between their middle objects carries both
frames to one another; one takes the equivalence relation generated by such
morphisms, makes the construction linear in the frames, and uses direct sum
for addition.  The resulting rational vector space is denoted \(A(X,Y)\).
We write \(\operatorname{Irr}(\mathcal P_i)\) for representatives of the
simple objects of \(\mathcal P_i\).
\end{definition}

The following length-three statement is the precise result behind the
phrase \emph{primitive equals an extension}.

\begin{lemma}[Three-level framed-object lemma]
\label{cp:lem-three-level-framed}
Let \(\mathcal C\) be an essentially small rational finite-length abelian
category with finite-dimensional Hom spaces and a
functorial strict decreasing filtration
\[
 M=F^0M\supset F^1M\supset F^2M\supset F^3M=0,
 \qquad \operatorname{gr}_F^iM\in\mathcal P_i,
\]
where the \(\mathcal P_i\) are semisimple.  Suppose each \(\mathcal P_i\)
has an exact faithful functor to finite-dimensional rational vector spaces.
For simple \(X\in\mathcal P_0\) and \(Y\in\mathcal P_2\),
\begin{equation}
\ker\!\left(
 \overline\Delta_{X,Y}:A(X,Y)\longrightarrow
 \bigoplus_{S\in\operatorname{Irr}(\mathcal P_1)}
 A(X,S)\otimes_{\operatorname{End}(S)}A(S,Y)
\right)
\cong\operatorname{Ext}^1_{\mathcal C}(X,Y).
\label{cp:eq-three-level-kernel}
\end{equation}
An extension \(0\to Y\to E\to X\to0\) corresponds to \(E\) framed by
the identity maps on its two graded pieces.
\end{lemma}

\begin{proof}
Strictness makes the associated-graded functor exact.  It is faithful as
well: if every graded map of a morphism is zero, strictness says that its
image has all graded pieces zero, hence the image itself is zero.  Composing
with the chosen fibre functors gives an exact faithful functor
\[
 \omega=\bigoplus_{i=0}^2\omega_i\operatorname{gr}_F^i:
 \mathcal C\longrightarrow\operatorname{Vect}^{\mathrm{fd}}_\Q.
\]

We recall the coefficient-coalgebra construction.  Start with
\[
 \bigoplus_{M\in\mathcal C}\omega(M)^\vee\otimes\omega(M)
\]
and impose, for every \(g:M\to N\), the matrix-coefficient relation
\[
 [\,\omega(g)^\vee\lambda\otimes x\,]_M
 =
 [\,\lambda\otimes\omega(g)x\,]_N.
\]
The quotient is a coalgebra: if \((e_a)\) and \((e_a^\vee)\) are dual
bases, its coproduct is
\[
 [\lambda\otimes x]_M\longmapsto
 \sum_a[\lambda\otimes e_a]_M\otimes[e_a^\vee\otimes x]_M.
\]
This quotient is also called the \emph{coend} of \(\omega\).  A coalgebra is
a vector space with a coassociative map to its tensor square; a comodule is
a vector space carrying a compatible coaction.  The
coefficient-coalgebra reconstruction theorem identifies
\(\mathcal C\) with its category of finite-dimensional comodules; it
requires an exact faithful fibre functor, but no tensor product on
\(\mathcal C\).  The spaces in
Definition~\ref{cp:def-coefficient-space} are its source--target
matrix-coefficient spaces: \(A(X,Y)\) is the block cut out by the idempotents
belonging to the pure simples \(X\) and \(Y\).

Because the pure categories are semisimple, their coefficient coalgebra is
cosemisimple, meaning that its finite-dimensional comodules are semisimple.
The normalized cobar resolution is the standard complex obtained by
repeatedly applying the coproduct after deleting the two pure endpoint
terms; it computes Yoneda extensions.  Thus the \emph{reduced} coproduct
below retains exactly the factorizations through the middle level.  In
source \(X\), target \(Y\), and
total level two, its first two positive-degree terms are
\[
 A(X,Y)\xrightarrow{\ \overline\Delta_{X,Y}\ }
 \bigoplus_{S\in\operatorname{Irr}(\mathcal P_1)}
 A(X,S)\otimes_{\operatorname{End}(S)}A(S,Y).
\]
There is no incoming degree-zero term because
\(\operatorname{Hom}_{\oplus\mathcal P_i}(X,Y)=0\) for distinct pure
levels.  To see the displayed differential directly, choose on every
\(S\)-isotypic part of \(\operatorname{gr}_F^1M\) dual maps
\[
 u_a:S\longrightarrow\operatorname{gr}_F^1M,\qquad
 u^a:\operatorname{gr}_F^1M\longrightarrow S.
\]
Then
\begin{equation}
 \overline\Delta_{X,Y}[M,v,f]
 =
 \sum_{S,a}
 [M/F^2M,v,u^a]\otimes_{\operatorname{End}(S)}
 [F^1M,u_a,f].
 \label{cp:eq-explicit-middle-coproduct}
\end{equation}
Thus degree-one cobar cohomology is the kernel in
\eqref{cp:eq-three-level-kernel}; reconstruction identifies it with
\(\operatorname{Ext}^1_{\mathcal C}(X,Y)\).
\end{proof}

For the category used here, take \(X=\Q_U(0)\) and \(Y=\Q_U(2)\).
The pure fibre functors are the usual forgetful functors on rational vector
spaces and on the rational Artin representations generated by \(\chi\).
Once the vanishings giving the strict filtration have been verified,
Lemma~\ref{cp:lem-three-level-framed} yields
\begin{equation}
 \ker\overline\Delta_{0,2}
 \cong
 \operatorname{Ext}^1_{\mathcal M_U^{[0,2]}}
   \bigl(\Q_U(0),\Q_U(2)\bigr)
 \cong H^1(U,\Q(2)).
 \label{cp:eq-primitive-is-ext}
\end{equation}
A framed class in this kernel is called \emph{primitive}.  In plain
language, once every factorization through the middle level has cancelled,
the only remaining obstruction is a direct extension of the top quotient
by the bottom subobject.

For a quadric tetrahedron, the six line sections \(D_I\) give the six
summands in \eqref{cp:eq-explicit-middle-coproduct}.  After a splitting
cover their two half-frames are weight-one cross-ratios: the angle
coordinate and the complementary edge coordinate.  Section
\ref{cp:sec-relative-rigidity} computes their sum.

For a field \(F\), weight-two motivic cohomology has the \(K\)-theoretic
description
\begin{equation}
 H^1(F,\Q(2))
 \cong \operatorname{gr}^{\gamma}_2K_3(F)\otimes\Q.
 \label{cp:eq-weight-two-k3}
\end{equation}
Here \(K_3(F)\) is Quillen's third algebraic \(K\)-group and
\(\operatorname{gr}^{\gamma}_2\) is the second associated-graded piece of
its standard gamma filtration.  No other part of algebraic \(K\)-theory is
used in the proof.
This explains why \(K_3\) appears at the end of a scissors-congruence
argument: it is the group classifying the primitive endpoint extensions in
Tate degree two.

\subsection{Specialization is pullback}
\label{cp:subsec-specialization}

Suppose the family \((Q,M,\alpha)\) is defined over \(U\), and let
\(s:\operatorname{Spec}\C\to U\) be a complex point.  Its specialization is
not an appeal to analytic continuity.  It is the exact functor \(s^*\).
Because every face term is constructible and dualizable, pullback gives
\[
 s^*\mathfrak m_U(Q,M)
 \cong \mathfrak m_{\C}(Q_s,M_s),
\]
and preserves the two frames, the three-step filtration, its coproduct, and
the endpoint triangle.  In particular, a zero relative framed class pulls
back to a zero framed class at every complex point.

If a primitive class is represented by
\[
 \Q_U(2)\longrightarrow E\longrightarrow\Q_U(0)
 \xrightarrow{\kappa_U}\Q_U(2)[1],
\]
then its fibre is represented by the literally pulled-back triangle and
endpoint \(s^*\kappa_U\).  This observation does not require a motivic
\(t\)-structure over \(\C\), nor does it require pullback to be \(t\)-exact:
the endpoint morphism already exists in the triangulated category before
pullback.

Goncharov's construction assigns to the specialized oriented quadric simplex
the same normalized framed face motive.  Hence its normalized-cobar endpoint
is exactly the canonical class in
\[
 H^1(\C,\Q(2))
 \cong \operatorname{gr}^{\gamma}_2K_3(\C)\otimes\Q,
\]
not merely a period or regulator of that class.

\subsection{Standard inputs and the boundary of the argument}
\label{cp:subsec-standard-inputs}

For clarity, the proof treats the following results as standard external
inputs.

\begin{enumerate}[label=(\arabic*)]
\item The rational constructible motivic formalism over the smooth bases in
      question: localization, purity, the projective bundle formula,
      internal duality, finite-\'etale trace, and exact symmetric monoidal
      pullback.  We also use cancellation and the higher-Chow description
      of motivic cohomology, in particular the vanishing in negative Tate
      degree and \(H^1(W,\Q(0))=0\) for the smooth schemes \(W\) occurring
      here.

\item Goncharov's finite semisimple gluing construction
      \cite[Chapter~5, Section~1]{Goncharov1999}, applied only to the three
      pure categories in Equation~\eqref{cp:eq-pure-three-levels}.

\item Goncharov's relative quadric motive
      \cite[Chapter~5, Section~2]{Goncharov1999}, together with the
      weight-coefficient and coproduct computations in
      \cite[Theorem~4.8 and Proposition~4.9]{Goncharov1999}.  The twist in
      Equation~\eqref{cp:eq-normalized-relative-motive} agrees with the
      normalization in \cite[Definition~5.8]{Brown2013}.

\item The coefficient-coalgebra reconstruction theorem and its normalized
      cobar resolution \cite{BGSV1990,Goncharov1999}.  Their complete
      length-three specialization is proved in
      Lemma~\ref{cp:lem-three-level-framed}.

\item Garkusha's weight-two birational comparison and invariance under a
      purely transcendental extension
      \cite[Corollaries~3.10, 3.11, 3.13 and Theorem~3.14]{Garkusha2023}.
      For the rational open base used later, these give, in degrees
      \(m\leq1\),
      \[
       H^m(U,\Q(2))\cong H^m(\Q(t_1,\ldots,t_6),\Q(2))
       \cong H^m(\Q,\Q(2)).
      \]

\item Beilinson--Soul\'e vanishing in the required nonpositive degrees over
      the number field \(\Q\), and Borel's rank theorem
      \cite{Borel1977}, which gives
      \[
       H^1(\Q,\Q(2))
       =\operatorname{gr}^{\gamma}_2K_3(\Q)\otimes\Q=0.
      \]
\end{enumerate}

\begin{theorem}[Weight-two comparison]
\label{cp:thm-goncharov-inputs}
The canonical map defined in Goncharov's equations \textup{(4.6)--(4.7)}
\cite{Goncharov1999} is valid and injective in weight two by his
Theorem~4.13.  His Theorem~4.14 supplies the compatible map from spherical
scissors classes to quadric simplices; the injectivity used to return from a
vanishing quadric class to a vanishing reduced spherical scissors class is
stated immediately after the proof of that theorem.
\end{theorem}

The argument does \emph{not} assume Goncharov's Conjecture~4.15, which asks
for a stronger generator-level chain map from the full quadric complex to a
Bloch complex.  Nor does it infer an algebraic identity from equality of
regulators.  The Regge-specific work that follows has three tasks only:
construct the common rational base, compute the middle coproduct and prove
that it vanishes there, and identify the resulting fibre with the canonical
class to which Theorem~\ref{cp:thm-goncharov-inputs} applies.
